% ------------------------------------------------------------------------
% Template zur Erstellung von Abschlussarbeiten an der Professur für
% Digitalisierung, E-Business und Operations Management
% Erstellerin: Jella Pfeiffer
% Datum 08.10.2019
% Überarbeitet von Pascal Heßler am 13.08.2020
% ------------------------------------------------------------------------

% ------------------------------------------------------------------------
% Allgemeine Einstellungen
% ------------------------------------------------------------------------
\documentclass[
    12pt,
    paper=a4,
    headings=small,               % kleinere Überschriften
    twoside=false,                % einseitig, nur rechte Seiten
    listof=totoc,                 % Listen im Inhaltsverzeichnis aufnehmen
    bibliography=totoc,           % Literaturverzeichnis ins Inhltsvz. aufnehmen
    headsepline                   % Trennlinie unter Kopfzeile
    ]{scrbook}

\usepackage[paper=a4paper,left=30mm,right=30mm,top=40mm,bottom=40mm]{geometry}
\usepackage[utf8]{inputenc}       % Damit können Umlaute ganz normal geschrieben werden.
\usepackage[english]{babel}       % Verwende deutsche, bzw. amerikanische Silbentrennung
%\usepackage[ngerman]{babel}      % Deutsche Alternative
\usepackage{scrlayer-scrpage}     % Paket für Kopf und Fußzeile. Kommentar: "scrpage2" veraltet in "scrlayer-scrpage" geändert
\pagestyle{scrheadings}           % Kopzeilenseitenstil
\usepackage{setspace}             % Zeilenabstand imlpementiert den Befehl 
% ------------------------------------------------------------------------
% Literatur
% ------------------------------------------------------------------------
% Achtung zur Einbindung der Literatur wird hier Biber verwendet und nich der Standard BibTex. In Sublime bracu man nichts verändern.
% 1. Sicherstellen, dass Biber installiert ist ->  in Tex Live vorinstalliert
% 2. Stellt eure Schreibumgebung um -> In TeXstudio: Optionen -> TeXstudio konfigurieren -> Erzeugen -> Standard Bibliographieprogramm
\usepackage[backend=biber,style=apa,natbib=true]{biblatex}
\DeclareLanguageMapping{american}{american-apa}         % Sprach Einstellung
%\DeclareLanguageMapping{ngerman}{ngerman-apa}          % Deutsche Alternative
\addbibresource{BibliographieAbschlussarbeit.bib}       % Legt die Datei fest in der die Literatur liegt (z. B.: Export aus Citavi)
\usepackage{breakcites}                                 % Falls Zitationen nicht ans Zeilenende passen
\usepackage{csquotes}                                   % Erweitert die Möglichkeiten beim zitieren
% ------------------------------------------------------------------------
% Mathematische Symbole
% ------------------------------------------------------------------------
\usepackage{amssymb}
\usepackage{amsmath}        % Formattierung von Tabellen und Matritzen
\usepackage{amsfonts}

% ------------------------------------------------------------------------
% Grafiken
% ------------------------------------------------------------------------
\usepackage{graphicx}       % Zum Einbinden von Grafiken
\usepackage{subfigure}      % Um Bilder innerhalb einer Figure anzuordnen (Bild a b c d...)

% ------------------------------------------------------------------------
% Tabellen
% Um einfache Tabellen zu erstellen kann https://www.tablesgenerator.com/ verwendet werden.
% ------------------------------------------------------------------------
\usepackage{tabularx}       % Ermöglicht weitere Tabellen Einstellungen
\usepackage{multirow}       % Ermöglicht Zellen zu verbinden
\usepackage{booktabs}       % Tabellen Midline Topline etc.
\usepackage{float}          % Tabelen Positionierung
\usepackage{longtable}      % Ermöglicht Tabellen über mehrere Seiten hinweg, so dass sie noch dasselbe Format besitzten ftp://ftp.fu-berlin.de/tex/CTAN/macros/latex/required/tools/longtable.pdf
% **********************************************************************************************************************
% Folgende Pakete sind rein optional und etwas komplexer in ihrer Umsetzung nur für etwas erfahrenere LaTeX User
% **********************************************************************************************************************
% \usepackage{siunitx}      % Ausrichten von Tabellen an dezimalstellen (Umsetzung etwas komplizierter)
% \usepackage{rccol}        % Ein Neues Format R (wird wahrscheinlich mit dem Folgenden R zu Problemen führen) und sorgt, dass nach , ausgerichtet wird und automatisches runden in Tabellen wird ermöglicht
% \usepackage{fltpoint}     % Gehört zu rccol
% \usepackage{dcolumn}      % Gehört zu rccol
%
% \newcolumntype{L}[1]{>{\raggedright\arraybackslash}p{#1}}       % Linksbündig mit Breitenangabe
% \newcolumntype{C}[1]{>{\centering\arraybackslash}p{#1}}         % Zentriert mit Breitenangabe
% \newcolumntype{R}[1]{>{\raggedleft\arraybackslash}p{#1}}        % Rechtsbündig mit Breitenangabe
% \newcolumntype{x}[1]{!{\centering\arraybackslash\vrule width #1}}

% ------------------------------------------------------------------------
% Sonstige
% ------------------------------------------------------------------------
\usepackage{scrhack}        % Löst unnötige Warnungen!
\usepackage{color}          % Ermöglicht Text einzufärben \pagecolor{FARBE}  \color{FARBE} \textcolor{FARBE}{TEXT} \colorbox{FARBE}{TEXT}
\usepackage{hyperref}       % Ermöglicht URLS schreibweise ist z.B. \url{http://www.uni-giessen.de}

\usepackage{eurosym}        % Für Euro Symbol \euro{} oder \eruo (Unterschieden sich bezüglich Leerzeichen danach)

\usepackage{framed}         % Einrahmen von Texten mit der shaded-Umgebung.

\usepackage{pdflscape}      % Querformat möglich über \begin{landscape} \end{landscape}
% ------------------------------------------------------------------------
% Fein Justierungen
% ------------------------------------------------------------------------
\DeclareMathOperator*{\argmax}{arg\,max}

\renewcommand{\floatpagefraction}{.9}     % vorher: .5
\newenvironment{packed_enum}{
\begin{enumerate}
  \setlength{\itemsep}{1pt}
  \setlength{\parskip}{0pt}
  \setlength{\parsep}{0pt}
}{\end{enumerate}}

%\setlength{\parindent}{0pt}                    % Kein Einzug beim Absatzbegin
\setlength{\parskip}{1pt plus 0pt minus 1pt}    % Abstand zwischen 2 Absätzen


% Kapitel Überschriften in Schriftart mit Serifen
\setkomafont{sectioning}{\normalfont\normalcolor\bfseries}
% Gestaltung der Kopfzeilen
\ohead{\pagemark} \cfoot{} \cohead{} \ihead{\headmark}
\setkomafont{pagehead}{\normalfont\bfseries}
%\setkomafont{pagenumber}{\normalfont\bfseries} \automark{section}

\hypersetup{hidelinks}                 % Versteckt die Links im PDF (optisch)
% ----- ende der präambel ----------------------------------
% Start des Dokuments
% ----------------------------------------------------------
\begin{document} 
\frontmatter                           % Corspann, Kapitel römisch nummeriert

%!TEX root = ../master.tex
% Die Titelseite der Arbeit

\begin{titlepage}

\begin{center} % zentrieren

  % Logo der Uni Giessen bitte suchen und einfuegen
  \begin{figure}[ht]
    \centering
    \includegraphics{graphics/logo.png}
  \end{figure}

  % Vertikaler Zwischenraum
  \bigskip
  \vfill
  % Titel der Arbeit und Typ der Arbeit, umrandet
    \begin{framed}
    \begin{center}
      \textsc{{\Large Titel\\
      Untertitel\\}}
                                % Letztes \\ ist wichtig, beginnt eine neue Zeile f{\"u}r die Art der Arbeit

      \bigskip

                                % Art der Arbeit, ggf. auszutauschen gegen Seminar- oder Bachelorarbeit
      \textbf{Masterthesis}
    \end{center}
    \end{framed}
    \vfill
    \vfill

%@Lukas Münzel: bitte austauschen durch eingereicht am, Name, Vorname, Matrikelnummer, Studiengang, private Adresse, Telefonnummer und E-Mail Adresse

  % Daten des Erstellers, Einreichungsdatum
  % in einer Tabelle ausgerichtet
  \begin{tabular*}{0.62\textwidth}{r@{\extracolsep{\fill}}l}
    eingereicht am: & DD.MM.YYYY \\\\
    von: & Vor- Nachname\\
    & geboren am DD.MM.YYYY \\
    \\
    Matrikelnummer: & XXX\\
    Studiengang: & XXX\\
    Private Adresse: & XXX\\
    Telefonnummer: & XXX\\
    E-Mail-Adresse: & XXX\\
    Betreuer: & XXX
  \end{tabular*}
  \vfill
  \vfill


  \rule{\textwidth}{.4pt}\\ % vertikale Linie
  Justus-Liebig-Universität Gießen\\
  Professur für Digitalisierung, E-Business und Operations Management\\
  Licherstraße 74 \\
  35394  Gießen\\
\end{center}

\end{titlepage} % Ende des Titelblatts

               % Titelseite einbinden
\tableofcontents                       % Fügt Inhaltsverzeichnis ein
% %!TEX root = ../master.tex
%-------------------------------------
%
%    Stichwortverzeichnis (ca. 5-10 Stichworte, welche den Inhalt der Arbeit beschreiben
%
%-------------------------------------

\chapter{Keywords} % beachte addchap
\begin{labeling}{1234567890}
        \item Duplicate Elimination
        \item Genetic Algorithm
        \item Greedy Heuristic
        \item Integer Program
        \item Meta-Heuristic
        \item Multi-Dimensional Knapsack Problem
        \item Multi-Unit Combinatorial Auction
        \item Phenotypic Distance
        \item Winner Determination Problem
\end{labeling}
            % Fügt Stichwortverzeichnis ein    (nur wenn nötig)
% \listoffigures                         % Fügt Abbildungsverzeichnis ein   (nur wenn nötig)
% \listoftables                          % Fügt Tabellenverzeichnis ein     (nur wenn nötig)
%!TEX root = ../master.tex
%-------------------------------------
%
%    minimales abk�rzungsverzeichnis
%
%-------------------------------------

\addchap{Abbreviations} % beachte addchap
\begin{labeling}{1234567890}
        \item[CA] combinatorial auction
        \item[VCG] Vickrey-Clarke-Grooves
\end{labeling}
       % Beispiel eines handerstellten Verzeichnisses Kommentar: Abkürzungsverzeichnisse können auch über \usepackage[acronyms, shortcuts, translate=babel, nomain]{glossaries} generiert werden.

% ------------------------------------------------------------------------
% Hauptarbeit (hier werden alle Kapitel eingebunden)
% ------------------------------------------------------------------------
\mainmatter                            % Hauptteil, Kapitel lateinisch nummeriert
\onehalfspacing
% Theoretisch kann man auch alles in einem Dokument haben. Das ist jedoch unübersichtlich. Daher werden hier die einzelnen Kapitel eingebunden.

\chapter{Fromeln}
Unter Verwenung der schnell schreibweise mittels eines einzelnen \emph{\$}-Zeichens kann relativ simpel eine Formel im Text verwendet werden.
$ c(\tau) =\frac{1}{|C|(|C|-1)} \sum_{j \in C\backslash \{i\}}1_{d_{ij}>\tau}$\\\\
Unter Verwenung der schnell schreibweise mittels zwei \emph{\$\$}-Zeichen wird die Formel in einem eigenen Absatz angezeigt:
$$ c(\tau) =\frac{1}{|C|(|C|-1)} \sum_{j \in C\backslash \{i\}}1_{d_{ij}>\tau}$$

Unter Verwenung der Umgebung \emph{align} wird die Formel automatisch Nummeriert.
\begin{align}
    c(\tau) =\frac{1}{|C|(|C|-1)} \sum_{j \in C\backslash \{i\}}1_{d_{ij}>\tau}
\end{align}

% ------------------------------------------------------------------------
% Hauptarbeit ENDE
% ------------------------------------------------------------------------

% ------------------------------------------------------------------------
% Erstellen des Literaturverzeichnisses
% ------------------------------------------------------------------------
\printbibliography[heading=bibintoc]  
% ------------------------------------------------------------------------
% Einbinden des Anhangs
% ------------------------------------------------------------------------
\appendix                               % Kapitel werden ab hier Alphabetisch aufgenommen
%!TEX root = ../master.tex
\chapter{Appendix}\label{Appendix}

\begin{table}[h]
\centering
\caption{Size of the search space for BASIC strategies ($n=4$).}
\label{tab:sizeSearchSpace}
 \begin{tabular}{r|r|r}
      m&s&search space size\\
      \hline
      4&1&64\\
      &2&4096\\
      &3&262,144\\
      &4&16,777,216\\
      &5&1,073,741,824\\
      7&1&262144\\
      &2&16777216\\
      &3&68,719,476,736\\
      &4&2.81E+14\\
      &5&1.15E+18\\
\end{tabular}
\end{table}



% ------------------------------------------------------------------------
% Einbinden der Eidesstaatlichen Erklärung
% ------------------------------------------------------------------------
\backmatter                             % Ab hier wird nichts mehr ins Inhaltsverzeichnis aufgenommen
%!TEX root = ../master.tex
\section*{Selbstständigkeitserklärung}
\thispagestyle{empty}
Hiermit versichere ich, die vorgelegte Thesis selbstständig und ohne unerlaubte fremde Hilfe und nur mit den Hilfen angefertigt zu haben, die ich in der Seminararbeit angegeben habe. Alle Textstellen, die wörtlich oder sinngemäß aus veröffentlichten Schriften entnommen sind, und alle Angaben die auf mündlichen Auskünften beruhen, sind als solche kenntlich gemacht. Bei den von mir durchgeführten und in der Thesis erwähnten Untersuchungen habe ich die Grundsätze guter wissenschaftlicher Praxis, wie sie in der \glqq Satzung der Justus-Liebig-Universität zur Sicherung guter wissenschaftlicher Praxisext\grqq{} niedergelegt sind, eingehalten. Gemäß \S 25 Abs. 6 der Allgemeinen Bestimmungen für modularisierte Studiengänge dulde ich eine Überprüfung der Thesis mittels Anti-Plagiatssoftware.
\bigskip

\raggedright{Gießen, den DD.MM.YYYY} \bigskip \bigskip \bigskip

Ihr NAME

\end{document}
% ----------------------------------------------------------
% Ende
% ----------------------------------------------------------
